%File: anonymous-submission-latex-2026.tex
\documentclass[letterpaper]{article} % DO NOT CHANGE THIS
\usepackage[submission]{aaai2027}  % DO NOT CHANGE THIS
\usepackage{times}  % DO NOT CHANGE THIS
\usepackage{helvet}  % DO NOT CHANGE THIS
\usepackage{courier}  % DO NOT CHANGE THIS
\usepackage[hyphens]{url}  % DO NOT CHANGE THIS
\usepackage{graphicx} % DO NOT CHANGE THIS
\urlstyle{rm} % DO NOT CHANGE THIS
\def\UrlFont{\rm}  % DO NOT CHANGE THIS
\usepackage{natbib}  % DO NOT CHANGE THIS AND DO NOT ADD ANY OPTIONS TO IT
\usepackage{caption} % DO NOT CHANGE THIS AND DO NOT ADD ANY OPTIONS TO IT
\frenchspacing  % DO NOT CHANGE THIS
\setlength{\pdfpagewidth}{8.5in} % DO NOT CHANGE THIS
\setlength{\pdfpageheight}{11in} % DO NOT CHANGE THIS
%
% These are recommended to typeset algorithms but not required. See the subsubsection on algorithms. Remove them if you don't have algorithms in your paper.
\usepackage{algorithm}
\usepackage{algorithmic}
\usepackage{graphicx} % Required for inserting images
\usepackage{amsmath,amssymb,amsthm}
\usepackage{amsfonts}
\usepackage{todonotes}
\usepackage{mathtools}

\theoremstyle{definition}
\newtheorem{lemma}{Lemma}
\theoremstyle{definition}
\newtheorem{definition}{Definition}
\theoremstyle{definition}
\newtheorem{theorem}{Theorem}
\theoremstyle{definition}
\newtheorem{proposition}{Proposition}

\theoremstyle{definition}

%
% These are are recommended to typeset listings but not required. See the subsubsection on listing. Remove this block if you don't have listings in your paper.
\usepackage{newfloat}
\usepackage{listings}
\DeclareCaptionStyle{ruled}{labelfont=normalfont,labelsep=colon,strut=off} % DO NOT CHANGE THIS
\lstset{%
	basicstyle={\footnotesize\ttfamily},% footnotesize acceptable for monospace
	numbers=left,numberstyle=\footnotesize,xleftmargin=2em,% show line numbers, remove this entire line if you don't want the numbers.
	aboveskip=0pt,belowskip=0pt,%
	showstringspaces=false,tabsize=2,breaklines=true}
\floatstyle{ruled}
\newfloat{listing}{tb}{lst}{}
\floatname{listing}{Listing}
%
% Keep the \pdfinfo as shown here. There's no need
% for you to add the /Title and /Author tags.
\pdfinfo{
/TemplateVersion (2026.1)
}

% DISALLOWED PACKAGES
% \usepackage{authblk} -- This package is specifically forbidden
% \usepackage{balance} -- This package is specifically forbidden
% \usepackage{color (if used in text)
% \usepackage{CJK} -- This package is specifically forbidden
% \usepackage{float} -- This package is specifically forbidden
% \usepackage{flushend} -- This package is specifically forbidden
% \usepackage{fontenc} -- This package is specifically forbidden
% \usepackage{fullpage} -- This package is specifically forbidden
% \usepackage{geometry} -- This package is specifically forbidden
% \usepackage{grffile} -- This package is specifically forbidden
% \usepackage{hyperref} -- This package is specifically forbidden
% \usepackage{navigator} -- This package is specifically forbidden
% (or any other package that embeds links such as navigator or hyperref)
% \indentfirst} -- This package is specifically forbidden
% \layout} -- This package is specifically forbidden
% \multicol} -- This package is specifically forbidden
% \nameref} -- This package is specifically forbidden
% \usepackage{savetrees} -- This package is specifically forbidden
% \usepackage{setspace} -- This package is specifically forbidden
% \usepackage{stfloats} -- This package is specifically forbidden
% \usepackage{tabu} -- This package is specifically forbidden
% \usepackage{titlesec} -- This package is specifically forbidden
% \usepackage{tocbibind} -- This package is specifically forbidden
% \usepackage{ulem} -- This package is specifically forbidden
% \usepackage{wrapfig} -- This package is specifically forbidden
% DISALLOWED COMMANDS
% \nocopyright -- Your paper will not be published if you use this command
% \addtolength -- This command may not be used
% \balance -- This command may not be used
% \baselinestretch -- Your paper will not be published if you use this command
% \clearpage -- No page breaks of any kind may be used for the final version of your paper
% \columnsep -- This command may not be used
% \newpage -- No page breaks of any kind may be used for the final version of your paper
% \pagebreak -- No page breaks of any kind may be used for the final version of your paperr
% \pagestyle -- This command may not be used
% \tiny -- This is not an acceptable font size.
% \vspace{- -- No negative value may be used in proximity of a caption, figure, table, section, subsection, subsubsection, or reference
% \vskip{- -- No negative value may be used to alter spacing above or below a caption, figure, table, section, subsection, subsubsection, or reference

\setcounter{secnumdepth}{0} %May be changed to 1 or 2 if section numbers are desired.

% The file aaai2026.sty is the style file for AAAI Press
% proceedings, working notes, and technical reports.
%

% Title

% Your title must be in mixed case, not sentence case.
% That means all verbs (including short verbs like be, is, using,and go),
% nouns, adverbs, adjectives should be capitalized, including both words in hyphenated terms, while
% articles, conjunctions, and prepositions are lower case unless they
% directly follow a colon or long dash
% \title{AAAI Press Anonymous Submission\\Instructions for Authors Using \LaTeX{}}
% \author{
%     %Authors
%     % All authors must be in the same font size and format.
%     Written by AAAI Press Staff\textsuperscript{\rm 1}\thanks{With help from the AAAI Publications Committee.}\\
%     AAAI Style Contributions by Pater Patel Schneider,
%     Sunil Issar,\\
%     J. Scott Penberthy,
%     George Ferguson,
%     Hans Guesgen,
%     Francisco Cruz\equalcontrib,
%     Marc Pujol-Gonzalez\equalcontrib
% }
% \affiliations{
%     %Afiliations
%     \textsuperscript{\rm 1}Association for the Advancement of Artificial Intelligence\\
%     % If you have multiple authors and multiple affiliations
%     % use superscripts in text and roman font to identify them.
%     % For example,

%     % Sunil Issar\textsuperscript{\rm 2},
%     % J. Scott Penberthy\textsuperscript{\rm 3},
%     % George Ferguson\textsuperscript{\rm 4},
%     % Hans Guesgen\textsuperscript{\rm 5}
%     % Note that the comma should be placed after the superscript

%     1101 Pennsylvania Ave, NW Suite 300\\
%     Washington, DC 20004 USA\\
%     % email address must be in roman text type, not monospace or sans serif
%     proceedings-questions@aaai.org
% %
% % See more examples next
% }

%Example, Multiple Authors, ->> remove \iffalse,\fi and place them surrounding AAAI title to use it
\title{From Weak Kleene to Classical Description Logics: \\Faithful Translations for Neuro-Symbolic Reasoning}
\author {{
    % Authors
    Thomas Ferguson\textsuperscript{\rm 1},
    James Oswald\textsuperscript{\rm 1},
}}
\affiliations {{
    % Affiliations
    \textsuperscript{\rm 1}Rensselaer Polytechnic Institute\\
    oswalj@rpi.edu, tferguson
}}

\newcommand{\strongneg}{\mathord{\sim}}
\newcommand{\fusion}{\mathrel{\circ}}
\newcommand{\isa}{\mathbin{{:}}}
\newcommand{\modi}{\mathcal{I}}
\newcommand{\vt}{\mathfrak{t}}
\newcommand{\vf}{\mathfrak{f}}
\newcommand{\vu}{\mathfrak{e}}
\newcommand{\vm}{\mathfrak{m}}
\newcommand{\vn}{\mathfrak{n}}
\newcommand{\vtt}{\langle\vt,\vt\rangle}
\newcommand{\vtf}{\langle\vt,\vf\rangle}
\newcommand{\vtu}{\langle\vt,\vu\rangle}
\newcommand{\vft}{\langle\vf,\vt\rangle}
\newcommand{\vff}{\langle\vf,\vf\rangle}
\newcommand{\vfu}{\langle\vf,\vu\rangle}
\newcommand{\vut}{\langle\vu,\vt\rangle}
\newcommand{\vuf}{\langle\vu,\vf\rangle}
\newcommand{\vuu}{\langle\vu,\vu\rangle}


\newtheorem{corollary}[theorem]{Corollary}


\begin{document}

\maketitle

\begin{abstract}

Recent research in the neuro-symbolic literature has empirically demonstrated that assessing large language model output with non-classical description logics (like $\mathcal{ALC}_\mathsf{AC}$ and $\mathcal{ALC}_{S^\star}$) leads to an increase in accuracy over traditional description logic methods. Left open, however, is the development of methods for the implementation of reasoning in these systems. In this work, we develop a finite translation scheme for embedding reasoning in $\mathcal{ALCH}_{\mathsf{AC}}$ within classical description logic with existential queries. We first translate weak Kleene truth and meaningfulness into a single classical knowledge base, linked by coherence axioms, and then translate the verification and falsification coordinates of $\mathsf{AC}$ into weak Kleene concepts. We prove that the composed translation preserves and reflects entailment, including arbitrary complex concepts, constants, and role hierarchies. The interpretation of existential queries is made explicit, and the proofs construct associated interpretations in both directions.

\end{abstract}

% Uncomment the following to link to your code, datasets, an extended version or similar.
% You must keep this block between (not within) the abstract and the main body of the paper.
% \begin{links}
%     \link{Code}{https://aaai.org/example/code}
%     \link{Datasets}{https://aaai.org/example/datasets}
%     \link{Extended version}{https://aaai.org/example/extended-version}
% \end{links}

\section{Introduction}

% \todo[inline]{
%     \begin{itemize}
%         \item What is ALCAC and ALCS\ast, where do they come from?
%         \item Why is reasoning in them desireable?
%         \begin{itemize}
%             \item Mention Fuzzy description logics.
%         \end{itemize}
%         \item Why do we need a translation?
%         \item Summary of what we do, main theorems, sections
%     \end{itemize}
% }

Description logics (DLs) provide a well-established formal foundation for representing structured knowledge. Their model-theoretic semantics, favorable computational properties, and use for ontology languages such as OWL have enabled the development of mature automated reasoning systems \cite{DBLP:conf/ore/MatentzogluLHSP15,DBLP:series/faia/SinghBM23}. Standard description logics, such as those described in \cite{DBLP:conf/dlog/2003handbook}, are classical: every assertion is evaluated relative to a two-valued interpretation, and a knowledge base is normally expected to admit a classical model. These assumptions are appropriate in many applications, but they can be too strong when the represented information is partial, semantically indeterminate, or internally conflicting.

Situations requiring the ability to handle non-classical reasoning like this arise naturally when a knowledge base represents information obtained from several agents or extracted from heterogeneous sources. For example, ontology-extraction systems, such as some of those described in \cite{DBLP:journals/csr/ArmaryVLN25}, may produce an assertion whose terminology is only partially aligned with that of an existing ontology. The assertion may conflict with the ontology without actually being false in the traditional sense: its failure may instead result from an undefined term, an unresolved mapping, or a lack of sufficient information. Similarly, when representing intentional states, an agent may accept one claim, reject another, and remain unable to evaluate a third. Treating all of these cases as ordinary classical falsity does not capture nuance that may be important for reasoning.

A number of non-classical description logics have been proposed to address limitations of classical two-valued semantics. One such example are fuzzy description logics \cite{DBLP:conf/aaai/Straccia98,DBLP:conf/sum/BorgwardtP17} which assign degrees of truth or membership to concepts and assertions. Fuzzy description logics are useful when the phenomenon being represented is continuous, such as similarity, vagueness, or uncertain classification. The problem considered here is different. We are interested not primarily in how strongly an assertion is true, but in whether it is truth-evaluable at all and in how positive and negative information concerning it interact. Consequently, the semantic distinctions required in our setting cannot in general be represented by a single numerical degree of truth.

The description logics $\mathcal{ALC}_\mathsf{AC}$ and $\mathcal{ALC}_{S^\ast}$, introduced as non-classical subsystems of $\mathcal{ALC}$ for modeling intentional states \cite{DBLP:conf/dlog/Ferguson21}, begin to address this problem. Their semantics are bilateral: formulas are evaluated using pairs of coordinates whose behavior is governed by weak Kleene operations. These coordinates permit different aspects of an evaluation, such as truth and meaningfulness, to be represented separately. These logics have recently acquired an additional motivation from neuro-symbolic reasoning with large language models. \cite{DBLP:conf/nesy/AllenCFIG25} propose grounding the interpretation function of a bilateral paraconsistent logic in a large language model, evaluating atomic formulas independently according to whether the model can verify and refute them. Their experiments show that bilateral factuality evaluation achieves higher macro-F1 than a unilateral factuality baseline. In domains in which factuality is critical---\cite{DBLP:conf/nesy/AllenCFIG25} considers examples from medical diagnosis---bilateral-based inference is a demonstrably useful tool.

The work in \cite{DBLP:conf/nesy/AllenCFIG25}, however, stops at the atomic case and does not consider sentences with the full expressivity offered by description logics. Moving from mere \emph{extraction} of reliable data from large language models to \emph{reliably reasoning} about such data is an important component in this plan. Increasing the utility of description logics like $\mathcal{ALC}_\mathsf{AC}$ and $\mathcal{ALC}_{S^\ast}$ thus becomes a pressing matter. 


%\todo[inline]{hard sell the connection here}
% Weak Kleene semantics are used since we take undefinedness to be infectious. If the evaluation of a compound expression depends essentially on a semantically undefined component, the compound expression may itself remain undefined rather than being assigned an arbitrary classical value. 





The additional expressive power in $\mathcal{ALC}_\mathsf{AC}$ and $\mathcal{ALC}_{S^\ast}$ creates a practical difficulty in terms of reasoning. Classical DLs benefit from several decades of work on reasoning procedures, optimization techniques, and highly developed software implementations. By contrast, specialized reasoners for $\mathcal{ALC}_\mathsf{AC}$ and $\mathcal{ALC}_{S^\ast}$ are not generally available. Implementing a new reasoning architecture directly from the non-classical semantics would require reproducing much of the infrastructure already present in classical description-logic reasoners. It would also make it more difficult to transfer advances in classical reasoning to the non-classical setting. We therefore investigate whether reasoning in these logics can instead be reduced to reasoning in a classical description logic. 

% \section{A Motivating Example}

% \todo[inline]{
%     Something similar to the example from \cite{DBLP:conf/dlog/Ferguson21}, but new.

%     Motivating example from CRAFT, when trying to extract and match ontologies, want to be able the degree to which something conflicts with a KB. 

%     A way of getting an assessment of weather a semantic extraction is inconsistent with a KB is to examine weaker notions of consistency with weaker description logics.
% }

\section{A Motivating Example}

Consider the following sample $\mathcal{ALCH}$ knowledge base about a person applying for a loan. Let $a$ be the loan applicant, $m$ be a mortgage account, and $c$ be a credit account.
\begin{align*}
R_1 ) \quad& \mathtt{hasCreditCardDebt} \sqsubseteq \mathtt{hasDebt}  \\
R_2) \quad& \mathtt{hasMortgageDebt} \sqsubseteq \mathtt{hasDebt} \\
A_1) \quad&\langle a,c \rangle : \mathtt{hasCreditCardDebt} \\
A_2) \quad&\langle a,m \rangle : \mathtt{hasMortgageDebt} \\
A_3) \quad&a:\forall \mathtt{hasDebt}.\mathtt{GoodStandingAcc} \\
A_4) \quad&m:\mathtt{GoodStandingAcc} \sqcap \mathtt{SecureAcc} \\
A_5) \quad&c : \lnot \mathtt{GoodStandingAcc}
\end{align*}
Statements in the knowledge-base may have been aggregated from multiple sources by a language model, leading to potential inconsistency. For example, $R_1$ and $A_1$ allow us to infer $\langle a,c\rangle : \mathtt{hasDebt} $ and together with $A_3$ this allows us to infer $c : \mathtt{GoodStandingAcc}$ which contradicts $A_5$. Under classical $\mathcal{ALCH}$ this knowledge base has no model and thus entails every formula including things like $m : \lnot \mathtt{SecuredAcc}$ or $a : \mathtt{FraudRisk}$ despite the knowledge base not containing evidence for either of these conclusions. By instead using paraconsistent $\mathcal{ALCH}$ variants we can represent that there is both evidence for and against the credit card account's good standing without allowing that conflict to spread arbitrarily.   


% Description logics such as $\mathcal{ALC}_{\mathsf{AC}}$ were introduced in the context of closing knowledge bases representing agents' \emph{intentional states} (\emph{e.g.}, beliefs or goals) under an appropriately weak notion of consequence. As agentic paradigms in artificial intelligence grow more entrenched, we continue to 

% A further motivating example is that of assessing the relative quality of semantic extraction from unstructured data. As Ma \emph{et al.} point out in \cite{DBLP:conf/dlog/MaQHL07},

% What we find to be the most pressing motivation follows from the recent research in neurosymbolic reasoning in Allen \emph{et al.}'s \cite{DBLP:conf/nesy/AllenCFIG25,allenetal2026}

% \todo[inline]{
% \begin{itemize}
%     \item For matching OWL, role subsumption is important so we're actually in ALC


%     $\mathcal{}$
% \end{itemize}

% }

\section{Background}

Here, we lay the foundations for our current work, by covering some preliminaries related to description logics and their non-classical extensions.


\subsection{Classical Description Logic}

We assume some familiarity with description logics, but give a brief coverage here, for a full account, see \cite{DBLP:conf/dlog/2003handbook}. Our translation will be presented from paraconsistent versions of $\mathcal{ALCH}$ to classical $\mathcal{ALCH^+(\cap)}$, but we note that the technique also works for paraconsistent $\mathcal{ALC}$ to classical $\mathcal{ALC^+}$, thus we present classical description logics up to $\mathcal{ALCH(\cap)}$ in a language independent way.
\begin{definition}[$\mathcal{ALC}$ Syntax] 
Let $\Sigma=(N_C,N_R,N_I)$ be a description logic signature, where $N_C$ is a set of atomic concept names, $N_R$ is a set of role names, and $N_I$ is a set of individual names. The set of $\mathcal{ALC}$ concepts over $\Sigma$ is generated by the grammar:
\[
C,D ::= \top \mid \bot \mid A \mid \lnot C \mid C\sqcap D
\mid C\sqcup D \mid \exists r.C \mid \forall r.C
\]
where $A\in N_C$ and $r\in N_R$. An $\mathcal{ALC}$ formula is defined by the following grammar where $a, b \in N_I$, $r \in N_R$ and $C, D$ are concepts. 
\[
\varphi ::= a \isa C \mid \langle a, b \rangle \isa r \mid C \sqsubseteq D \mid C \equiv D
\] 
\end{definition}
\noindent Sentences 1) of the form $a \isa C$ are called {\it concept assertions}, read ``$a$ is a $C$'',
2) of the form $(a, b) \isa r$ are called {\it role assertions}, read ``$a$ bears $r$ to $b$'', 3) of the form $C \sqsubseteq D$ are called {\it concept subsumptions}, read ``All $C$s are $D$s'', 4) of the form $C \equiv D$ are called {\it concept equivalences} read ``$C$ is equivalent to $D$''.

\begin{definition}[$\mathcal{ALCH}$, $\mathcal{ALC}^+$ $\mathcal{ALCH^+}$, $\mathcal{ALCH^+(\cap)}$]
Extending the sentence gammer of $\mathcal{ALC}$ with role inclusions $r \sqsubseteq s$ yields the language $\mathcal{ALCH}$. Sentences of this form can be read ``$r$ is a sub-role of $s$''. Removing the negation concept constructor $\lnot C$ from the concept gammer yields "positive version" of the language which we denote $\cdot^+$ (for readers familiar with DLs, note that $\mathcal{ALC^+}$ is not the same as the more restrictive $\mathcal{AL}$). Adding a role language $R,S::= A |R \sqcap S$ where $A \in N_R$ with a single constructor $R \cap S$ to any of these languages yields a $(\cap)$ version of that language, where roles can be joined, i.e. $(a, b) \isa r \cap s$ can be read ``$a$ bears $r$ to $b$ and $a$ bears $r$ and $s$ to $b$''. We note that for $(\cap)$ languages, we do not allow role intersection in concept restrictions.
\end{definition}

\begin{definition}[Knowledge Bases]
    Given a language $L$ an $L$ \textit{knowledge base} $\mathcal{K} = \langle\mathcal{T}, \mathcal{A}\rangle$ is a pair composed of 1) An \textit{ABox}, $\mathcal{A}$, a finite set of $L$ sentences that are concept assertions and role assertions and 2) A {\it TBox}, $\mathcal{T}$, a finite set of $L$ sentences that are concept subsumptions, concept equivalences and role subsumptions in $\mathcal{H}$ languages.
\end{definition}

% We present now present the classical semantics in terms of the most expressive language we use in this work $\mathcal{ALCH}(\cap)$, less expressive variants can be obtained by just removing any interpretation conditions that contain a feature not present in the given language.  

\begin{definition}[Classical Interpretations]\label{def:ALC-Interp-Non-Standard}\footnote{This is a nonstandard representation of interpretations, using the characteristic function of standard interpretations readers familiar with \cite{DBLP:conf/dlog/2003handbook} will know. This representation is meant to better prepare the reader for our later non-classical interpretations in Definitions \ref{def:wK-Interp} and \ref{def:Bilateral-Interp}.} Let $\mathcal{V}_2 = \{\vt,\vf\}$ be the set of classical truth values and let $L$ be our language. An $L$-interpretation consists of a nonempty domain $\Delta^{\mathcal{I}}$ and interpretation function $\mathcal{I}$ from $L$ formulae to $\mathcal{V}_2$ such that the following hold depending on $L$: For $\mathcal{ALC}$ Languages:
\begin{itemize}
		\item $\mathcal{I}(a\isa\top)=\vt$ and $\mathcal{I}(a\isa\bot)=\vf$
		\item $\mathcal{I}(a\isa C)\in\mathcal{V}_{2}$ and $\mathcal{I}(\langle a,b\rangle\isa r)\in\mathcal{V}_{2}$
        \item $\mathcal{I}(a\isa\neg C)=\lnot\mathcal{I}(a\isa C)$ (not in $\cdot^+$ languages)
		\item $\mathcal{I}(a\isa C\sqcap D)=\mathcal{I}(a\isa C)\land\mathcal{I}(a\isa D)$
		\item $\mathcal{I}(a\isa C\sqcup D)=\mathcal{I}(a\isa C)\lor\mathcal{I}(a\isa D)$
		\item $\mathcal{I}(a\isa\exists r.C) =\exists {b\in\Delta^{\mathcal{I}}}, \left( \mathcal{I}(\langle a,b\rangle\isa r) \land \mathcal{I}(b\isa C) \right)$ \item $\mathcal{I}(a\isa\forall r.C) =\forall {b\in\Delta^{\mathcal{I}}}, \left( \mathcal{I}(\langle a,b\rangle\isa r) \rightarrow \mathcal{I}(b\isa C) \right)$ \item $\mathcal{I}(C\sqsubseteq D) =\forall{a\in\Delta^{\mathcal{I}}}, \left( \mathcal{I}(a\isa C) \rightarrow \mathcal{I}(a\isa D) \right)$ \item $\mathcal{I}(C\equiv D) =\mathcal{I}(C\sqsubseteq D) \land \mathcal{I}(D\sqsubseteq C)$
    \end{itemize}
    \noindent For $\mathcal{H}$ languages:
    \begin{itemize}
        \item $\mathcal{I}(r \sqsubseteq s) = \forall a,b \in \Delta^\mathcal{I}, \mathcal{I}(\langle a, b\rangle : r) \rightarrow \mathcal{I}(\langle a, b\rangle : s)$
	\end{itemize}
    \noindent For $(\cap)$ languages:
    \begin{itemize}
        \item $\mathcal{I}(\langle a, b\rangle : r \sqcap s) = \mathcal{I}(\langle a, b\rangle : r) \land \mathcal{I}(\langle a, b\rangle : s)$
    \end{itemize}
\end{definition}

\begin{definition}[Satisfaction]\label{def:ALC-Sat}
    Given an $L$-interpretation $\mathcal{I}$ and formula $\varphi$ in $L$, we say $\mathcal{I}$ satisfies $\varphi$, or conversely $\varphi$ holds in $\mathcal{I}$, written $\mathcal{I}\models_L\varphi$, iff $\mathcal{I}(\varphi) = \vt$. With this, we can say a $L$ knowledge base $\mathcal{K} = \langle \mathcal{T},\mathcal{A}\rangle$ holds under an $L$-interpretation $\mathcal{I}$, written $\mathcal{I}\models_L\mathcal{K}$ iff $\forall \varphi \in \mathcal{A}, \mathcal{I}\models_L\varphi \text{ and } \forall \varphi \in \mathcal{T}, \mathcal{I}\models_L\varphi$. 
\end{definition}

\noindent 
\begin{definition}[Entailment]\label{def:ALC-Ent}
We say a $L$ knowledge-base $\mathcal{K}$ entails a formula $\varphi$, written $\mathcal{K} \models_{L} $ iff for every interpretation $\mathcal{I}$, $\mathcal{I}\models_L\mathcal{K}$ implies $\mathcal{I}\models_L\varphi$.
\end{definition}

\subsection{Weak Kleene, $\mathsf{AC}$, and $\mathsf{S}^*$}
To understand $\mathcal{ALCH}_\mathsf{AC}$ and $\mathcal{ALCH}_{\mathsf{S}^*}$ we must first discuss the underlying logics of \textit{anaclitic containment}, $\mathsf{AC}$, introduced in \cite{Angell1977-ANGTSO-9} and $\mathsf{S}^*$, introduced in \cite{Daniels1990-DANANO-3}, as well as their construction from weak Kleene logic. The syntax of both are the same as classical propositional logic with two additional restricted quantifier connectives. The truth values of $\mathsf{AC}$ and $\mathsf{S}^*$ are based on the three truth values, $\mathcal{V}_3 = \{\vt, \vu, \vf\}$, from weak Kleene logic $\mathsf{wK}$. Our semantics will also depend on the $\mathsf{wK}$ connectives and quantifiers.  
\begin{definition}[$\mathsf{wK}$ Truth Tables]
\leavevmode\par\noindent
\begin{tabular}{c|c c c|ccc c c|ccc}
    $\dot{\strongneg}$ & && $\dot{\wedge}$ & $\vt$ & $\vu$ & $\vf$ && $\dot{\vee}$ & $\vt$ & $\vu$ & $\vf$\\
    \cline{1-2}\cline{4-7}\cline{9-12}
    $\vt$ & $\vf$ && $\vt$ & $\vt$ & $\vu$ & $\vf$ && $\vt$ & $\vt$ & $\vu$ & $\vt$\\
    $\vu$ & $\vu$ && $\vu$ & $\vu$ & $\vu$ & $\vu$ && $\vu$ & $\vu$ & $\vu$ & $\vu$\\
    $\vf$ & $\vt$ && $\vf$ & $\vf$ & $\vu$ & $\vf$ && $\vf$ & $\vt$ & $\vu$ & $\vf$\\
\end{tabular}
\end{definition}
\noindent Due to issues observed in \cite{DBLP:conf/tableaux/Ferguson21,DBLP:conf/dlog/Ferguson21} we use the following \textit{restricted Kleene quantifiers} that operate on sets of pairs of $\mathsf{wK}$ truth values. The motivation for these quantifiers is that restricted quantification statements such as "all $\varphi$s are $\psi$s" ($[\forall x\varphi(x)]\psi(x)$) are truth-evaluable only when the domains of $\varphi$ and $\psi$ can be compared.
% Local spacing keeps the unchanged original definition within the column.
\begingroup
\thickmuskip=3mu
\medmuskip=3mu
\begin{definition}(Restricted $\mathsf{wK}$ Quantifiers)
    The \emph{restricted Kleene quantifiers} are functions $\dot{\exists}$ and $\dot{\forall}$ mapping a nonempty sets $X\subseteq\mathcal{V}_{3}^{2}$ to truth values from $\mathcal{V}_{3}$ as follows:
	
	\vspace{.1cm}

    {\small
	
	\noindent $\dot{\exists}(X)=\begin{cases}
	\vt & \mbox{ if }\langle\vt,\vt\rangle\in X\\
	\vu & \mbox{ if }\forall\langle u,v\rangle\in X, u=\vu\lor v=\vu\\
	\vf & \mbox{ if }\langle\vt,\vt\rangle\notin X \land (\exists \langle u,v\rangle\in X\mbox{, }u\neq\vu\land v\neq\vu)
	\end{cases}$

    	\noindent $\dot{\forall}(X)=\begin{cases}
	\vt & \mbox{ if }\langle\vt,\vf\rangle,\langle\vt,\vu\rangle\notin X \land (\exists \langle u,v\rangle\in X\mbox{, }u\neq\vu\land v\neq\vu)\\
	\vu & \mbox{ if } \forall\langle u,v\rangle\in X,u=\vu\lor v=\vu\\
	\vf & \mbox{ if }
    \begin{aligned}
    \lbrace\langle\vt,\vf\rangle,\langle\vt,\vu\rangle\rbrace\cap X\neq\varnothing \mbox{ and }\\ \exists \langle u,v\rangle\in X\mbox{, }u\neq\vu\land v\neq\vu\\
    \end{aligned}
	\end{cases}$
    }
\end{definition}
\endgroup
\noindent We now turn to the construction of $\mathsf{AC}$ from $\mathsf{wK}$. $\mathsf{AC}$ is \textit{bilateral} in that for its truth values, it assigns statements pairs of $\mathsf{wK}$ truth values, $\mathcal{V}^2_3 \: (\mathcal{V}_3 \times \mathcal{V}_3)$ where the first coordinate represents the degree of support (evidence) for a statement, while the second represents the degree of refutation of (evidence against) a statement. This yields nine possible truth values in $\mathsf{AC}$. One interpretation of these truth values is given as follows:
\begin{align*} \langle\vt,-\rangle: & \text{ evidence exists for } \varphi \\ 
\langle -,\vt\rangle: & \text{ evidence exists refuting } \varphi \\ 
\langle\vu,-\rangle: & \text{ the truth conditions of } \varphi \text{ are undefined} \\
\langle -,\vu\rangle: & \text{ the falsity conditions of } \varphi \text{ are undefined} \\ 
\langle\vf,-\rangle: & \text{ no evidence exists for } \varphi \\
\langle -,\vf\rangle: & \text{ no evidence exists refuting } \varphi 
\end{align*}
With these truth values, we now turn to defining connectives and quantifiers for them.
\begin{definition}[$\mathsf{AC}$ Connectives]\label{def:AC-connectives}
The $\mathsf{AC}$ truth functions $\ddot{\strongneg}$, $\ddot{\wedge}$, and $\ddot{\vee}$ are defined:
	\begin{align*}
	    \mathop{\ddot{\strongneg}}(\langle u,v\rangle)&:=\langle v,u\rangle \\
		\langle u_{0},v_{0}\rangle\mathrel{\ddot{\wedge}}\langle u_{1},v_{1}\rangle&:=\langle u_{0}\mathrel{\dot{\wedge}}u_{1}, v_{0}\mathrel{\dot{\vee}}v_{1}\rangle \\
		\langle u_{0},v_{0}\rangle\mathrel{\ddot{\vee}}\langle u_{1},v_{1}\rangle&:=\langle u_{0}\mathrel{\dot{\vee}}u_{1}, v_{0}\mathrel{\dot{\wedge}}v_{1}\rangle 
	\end{align*}
\end{definition}
\begin{definition}[Restricted $\mathsf{AC}$ Quantifiers]\label{def:AC-quantifiers}
Restricted quantifiers for $\mathsf{AC}$ are defined over sets $X$ of pairs of values from $\mathcal{V}^{2}_{3}$. Let $\pi^{i,j}$ map values $\langle\langle u_{0},v_{0}\rangle,\langle u_{1},v_{1}\rangle\rangle$ to the pair of the $i$th coordinate of $\langle u_{0},v_{0}\rangle$ and the $j$th coordinate of $\langle u_{1},v_{1}\rangle$ and let $\pi^{i,j}[X]$ be the image of $X$ under $\pi^{i,j}$. Then:
The $\mathsf{AC}$ restricted quantifiers $\ddot{\exists}$ and $\ddot{\forall}$ are defined to work with:
	\begin{align*}
		\ddot{\exists}(X)&:=\langle \dot{\exists} (\pi^{0,0}[X]),\dot{\forall}(\pi^{0,1}[X])\rangle \\
		\ddot{\forall}(X)&: =\langle \dot{\forall} (\pi^{0,0}[X]),\dot{\exists}(\pi^{0,1}[X])\rangle
	\end{align*}
\end{definition}
\noindent Finally, we turn to $\mathsf{S}^*$. For truth values, $\mathsf{S}^*$ uses the five-valued subset $\{\langle \vt, \vt \rangle, \langle\vt,\vf\rangle,\langle\vu,\vu\rangle,\langle\vf,\vt\rangle,\langle \vf, \vf \rangle\}$ of $\mathcal{V}_3^{2}$, which we denote $\mathcal{V}_3^{2*}$. This set of truth values is given due to philosophical objections about the existence of truth values like $\langle \vt, \vu \rangle$ which suggest that one could be aware of the truth conditions for a statement without awareness of its falsity conditions, and vice-versa.\footnote{See \cite{b108fe03530e4bd0ab883a7a5c2b5cf2} for an argument something is meaningful precisely when its negation is meaningful.} We do not make a commitment either way, and present the rest of this work in terms of the availability of both. For connectives and restricted quantifiers, $\mathsf{S}^*$ uses the same ones presented for $\mathsf{AC}$ in Definitions \ref{def:AC-connectives} and \ref{def:AC-quantifiers} but with its restricted set of truth values $\mathcal{V}_3^{2*}$. We note that this works because $\mathcal{V}_3^{2*}$ is closed under the $\mathsf{AC}$ connectives and restricted quantifiers.

\subsection{The Logics $\mathcal{ALCH^+(\cap)}_\mathsf{wK}$, $\mathcal{ALCH}_\mathsf{AC}$, and $\mathcal{ALCH}_{\mathsf{S}^*}$}

We begin by presenting the description logic  $\mathcal{ALCH^+(\cap)}_\mathsf{wK}$ which we will use as an intermediate step in our translations. $\mathcal{ALCH^+(\cap)}_\mathsf{wK}$ uses the language of classical $\mathcal{ALCH(\cap)}$ but for semantics uses the following $\mathsf{wK}$ interpretations:

\begin{definition}[$\mathsf{wK}$ Interpretations]\label{def:wK-Interp}
An $\mathcal{ALCH^+(\cap)}_\mathsf{wK}$ interpretation is a function $\mathcal{I}$ that maps $\mathcal{ALCH^+(\cap)}$ sentences to $\mathcal{V}_3$, paired with domain $\Delta^{\mathcal{I}}$, such that:
	\begin{itemize}
		\item $\mathcal{I}(a:\top)=\mathfrak{t}$ and $\mathcal{I}(a : \bot) = \vf$ for all $a\in\Delta^{\mathcal{I}}$.
		\item $\mathcal{I}(a\isa C)\in\mathcal{V}_{3}$ and $\mathcal{I}(\langle a,b\rangle\isa r)\in\mathcal{V}_{3}$
		% \item $\mathcal{I}(a\isa\neg C)=\mathop{\dot{\strongneg}}\mathcal{I}(a\isa C)$
		\item $\mathcal{I}(a\isa C\sqcap D)=\mathcal{I}(a\isa C)\mathrel{\dot{\wedge}}\mathcal{I}(a\isa D)$
		\item $\mathcal{I}(a\isa C\sqcup D)=\mathcal{I}(a\isa C)\mathrel{\dot{\vee}}\mathcal{I}(a\isa D)$
		\item $\mathcal{I}(a\isa \exists r.C)=\mathop{\dot{\exists}}(\lbrace\langle\mathcal{I}(\langle a,b\rangle\isa r),\mathcal{I}(b\isa C)\rangle\mid b\in \Delta^{\mathcal{I}}\rbrace)$
		\item $\mathcal{I}(a\isa \forall r.C)=\mathop{\dot{\forall}}(\lbrace\langle\mathcal{I}(\langle a,b\rangle\isa r),\mathcal{I}(b\isa C)\rangle\mid b\in \Delta^{\mathcal{I}}\rbrace)$
		\item $\mathcal{I}(C\sqsubseteq D)=\mathop{\dot{\forall}}(\lbrace\langle\mathcal{I}(a\isa C),\mathcal{I}(a\isa D)\rangle\mid a\in \Delta^{\mathcal{I}}\rbrace)$
		\item $\mathcal{I}(C\equiv D)=\mathcal{I}(C\sqsubseteq D)\mathrel{\dot{\wedge}}\mathcal{I}(D\sqsubseteq C)$
        \item $\mathcal{I}(r \sqsubseteq s) = \mathop{\dot{\forall}}(\lbrace\langle\mathcal{I}(\langle a,b\rangle\isa r),\mathcal{I}(\langle a,b\rangle\isa s)\rangle\mid a,b\in \Delta^{\mathcal{I}}\rbrace)$
        \item $\mathcal{I}(\langle a, b\rangle : r \sqcap s) = \mathcal{I}(\langle a, b\rangle : r) \mathrel{\dot{\wedge}} \mathcal{I}(\langle a, b\rangle : s)$
	\end{itemize}
\end{definition}
\noindent From this, the definitions of \textsf{wK} satisfaction $\mathcal{I}\models_{\mathcal{ALCH^+(\cap)}_\mathsf{wK}}\varphi$ and entailment $\mathcal{K}\models_{\mathcal{ALCH^+(\cap)}_\mathsf{wK}}\varphi$ are defined identically to the classical versions Definitions \ref{def:ALC-Sat} and \ref{def:ALC-Ent} using $\mathsf{wK}$ Interpretations.
% \begin{definition}[$\mathcal{ALCH^+(\cap)}_\mathsf{wK}$ Satisfaction]\label{def:wK-Sat}
%     An $\mathcal{ALCH^+(\cap)}_\mathsf{wK}$ interpretation $\mathcal{I}$ satisfies a $\mathcal{ALCH(\cap)}$ formula $\varphi$, written $\mathcal{I}\models_{\mathcal{ALCH^+(\cap)}_\mathsf{wK}}\varphi$, iff $\mathcal{I}(\varphi) = \vt$. A knowledge base $\mathcal{K} = \langle \mathcal{T},\mathcal{A}\rangle$ holds under an $\mathsf{wK}$ interpretation $\mathcal{I}$, written $\mathcal{I}\models_{\mathcal{ALCH^+(\cap)}_\mathsf{wK}}\mathcal{K}$ iff $\forall \varphi \in \mathcal{A}, \mathcal{I}\models_{\mathcal{ALCH^+(\cap)}_\mathsf{wK}}\varphi \text{ and } \forall \varphi \in \mathcal{T}, \mathcal{I}\models_{\mathcal{ALCH^+(\cap)}_\mathsf{wK}}\varphi$. 
% \end{definition}
% \begin{definition}[$\mathcal{ALCH^+(\cap)}_\mathsf{wK}$ Entailment]\label{def:wK-Ent}
% A $\mathcal{ALCH^+(\cap)}_\mathsf{wK}$ knowledge-base $\mathcal{K}$ entails a $\mathcal{ALCH(\cap)}$ formula $\varphi$, written $\mathcal{K} \models_{\mathcal{ALCH^+(\cap)}_\mathsf{wK}} \varphi $ iff for every $\mathcal{ALCH^+(\cap)}_\mathsf{wK}$ interpretation $\mathcal{I}$, $\mathcal{I}\models_{\mathcal{ALCH^+(\cap)}_\mathsf{wK}}\mathcal{K}$ implies $\mathcal{I}\models_{\mathcal{ALCH^+(\cap)}_\mathsf{wK}}\varphi$.
% \end{definition}

With our intermediate language defined, we now turn to defining the description logics $\mathcal{ALCH}_\mathsf{AC}$ and $\mathcal{ALCH}_{\mathsf{S}^*}$. They are both defined using $\mathcal{ALCH}$ as their underlying language. For their semantics we make use of our AC connectives and quantifiers from Definitions \ref{def:AC-connectives} and \ref{def:AC-quantifiers}.

\begin{definition}[$\mathsf{AC}$ Interpretations]\label{def:Bilateral-Interp}
	An $\mathcal{ALCH}_\mathsf{AC}$ interpretation is a function $\mathcal{I}$ from $\mathcal{ALCH}$ sentences to $\mathcal{V}_{3}^{2}$ paired with domain $\Delta^{\mathcal{I}}$, such that:
	
	\begin{itemize}
		\item $\mathcal{I}(a\isa\top)=\langle\vt,v\rangle$ and $\mathcal{I}(a\isa\bot)=\langle v,\vt\rangle$ for some $v\in\mathcal{V}_{3}$
		\item $\mathcal{I}(a\isa C)\in\mathcal{V}_{3}^{2}$ and $\mathcal{I}(\langle a,b\rangle\isa r)\in\mathcal{V}_{3}^{2}$
		\item $\mathcal{I}(a\isa\neg C)=\mathop{\ddot{\strongneg}}\mathcal{I}(a\isa C)$
		\item $\mathcal{I}(a\isa C\sqcap D)=\mathcal{I}(a\isa C)\mathrel{\ddot{\wedge}}\mathcal{I}(a\isa D)$
		\item $\mathcal{I}(a\isa C\sqcup D)=\mathcal{I}(a\isa C)\mathrel{\ddot{\vee}}\mathcal{I}(a\isa D)$
		\item $\mathcal{I}(a\isa \exists r.C)=\mathop{\ddot{\exists}}(\lbrace\langle\mathcal{I}(\langle a,b\rangle\isa r),\mathcal{I}(b\isa C)\rangle\mid b\in \Delta^{\mathcal{I}}\rbrace)$
		\item $\mathcal{I}(a\isa \forall r.C)=\mathop{\ddot{\forall}}(\lbrace\langle\mathcal{I}(\langle a,b\rangle\isa r),\mathcal{I}(b\isa C)\rangle\mid b\in \Delta^{\mathcal{I}}\rbrace)$
		\item $\mathcal{I}(C\sqsubseteq D)=\mathop{\ddot{\forall}}(\lbrace\langle\mathcal{I}(a\isa C),\mathcal{I}(a\isa D)\rangle\mid a\in \Delta^{\mathcal{I}}\rbrace)$
		\item $\mathcal{I}(C\equiv D)=\mathcal{I}(C\sqsubseteq D)\mathrel{\ddot{\wedge}}\mathcal{I}(D\sqsubseteq C)$
        \item $\mathcal{I}(r\sqsubseteq s)=\ddot{\forall}(\lbrace \langle\mathcal{I}(\langle a,b\rangle:r),\mathcal{I}(\langle a,b\rangle:s)\rangle\mid a,b\in\Delta\rbrace)$
	\end{itemize}
\end{definition}

\begin{definition}[$\mathcal{ALCH}_\mathsf{AC}$ Satisfaction]\label{def:AC-Sat}
    An $\mathsf{AC}$ interpretation $\mathcal{I}$ satisfies a formula $\varphi$, written $\mathcal{I}\models_{\mathcal{ALCH}_\mathsf{AC}}\varphi$, iff the first coordinate of $\mathcal{I}(\varphi)$ is $\vt$. A knowledge base $\mathcal{K} = \langle \mathcal{T},\mathcal{A}\rangle$ holds under an $\mathsf{AC}$ interpretation $\mathcal{I}$, written $\mathcal{I}\models_{\mathcal{ALCH}_\mathsf{AC}}\mathcal{K}$ iff $\forall \varphi \in \mathcal{A}, \mathcal{I}\models_{\mathcal{ALCH}_\mathsf{AC}}\varphi \text{ and } \forall \varphi \in \mathcal{T}, \mathcal{I}\models_{\mathcal{ALCH}_\mathsf{AC}}\varphi$. 
\end{definition}
\begin{definition}[$\mathcal{ALCH}_\mathsf{AC}$ Entailment]\label{def:AC-Ent}
An $\mathcal{ALCH}_\mathsf{AC}$ knowledge-base $\mathcal{K}$ entails a formula $\varphi$, written $\mathcal{K} \models_{\mathcal{ALC}_{\mathsf{AC}}} \varphi$ iff for every $\mathsf{AC}$ interpretation $\mathcal{I}$, $\mathcal{I}\models_{\mathcal{ALCH}_\mathsf{AC}}\mathcal{K}$ implies $\mathcal{I}\models_{\mathcal{ALCH}_\mathsf{AC}}\varphi$.
\end{definition}
These $\mathsf{AC}$ definitions generalize to $\mathcal{ALCH}_{\mathsf{S}^*}$ definitions by restricting the truth values in $\mathcal{ALCH}_\mathsf{AC}$ interpretations to exclusively be in $\mathcal{V}_3^{2*}$. Now with our languages and their semantics defined, we turn to the core of this work, providing a faithful translation for converting entailment problems in $\mathcal{ALCH}_\mathsf{AC}$ knowledge bases to classical $\mathcal{ALCH^+(\cap)}$ with explicitly interpreted existential queries.  

\section{Translation Framework}

The primary contribution of this paper lies in our describing techniques for reasoning with the above description logics by means of a classical translation. We proceed in two stages: first translating $\mathcal{ALCH}^{+}(\cap)_{\mathsf{wK}}$ into a classical language, and then decomposing $\mathcal{ALCH}_{\mathsf{AC}}$ into weak Kleene coordinates. The classical interpretation conditions remain those of Definition~\ref{def:ALC-Interp-Non-Standard}. We additionally give an explicit existential meaning to the blank-node assertions used in the translation.

\subsection{Embedding $\mathcal{ALCH}^+(\cap)_{\mathsf{wK}}$ in $\mathcal{ALCH}^+(\cap)$}

We begin with our intermediate $\mathcal{ALCH}^+(\cap)_{\mathsf{wK}}$. For every atomic concept $C$ and atomic role $r$ in its signature, introduce distinct classical copies $C^{\tau},C^{\mu}$ and $r^{\tau},r^{\mu}$. For concepts, $\cdot^{\tau}$ records truth and $\cdot^{\mu}$ records meaningfulness. The copies must be interpreted together: truth implies meaningfulness, and the truth of a restricted universal statement requires a jointly meaningful witness.

The latter requirement prevents an atomwise replacement of meaningful values by $\vt$ from supplying a translation. For example, $\dot{\forall}(\{\langle\vf,\vf\rangle,\langle\vf,\vu\rangle\})=\vt$, whereas replacing the non-$\vu$ coordinates by $\vt$ gives $\dot{\forall}(\{\langle\vt,\vt\rangle,\langle\vt,\vu\rangle\})=\vf$. We instead preserve truth and meaningfulness simultaneously by structural induction.

\begin{definition}[Blank-node assertions]\label{def:blank-nodes}
In the classical target, $\_:C$ is satisfied when some element of the domain satisfies $C$. Likewise, $\langle\_,\_\rangle:R$ is satisfied when some pair of domain elements satisfies $R$. Thus:
\begin{align*}
\mathcal{I}\models\_:C
 &\quad\text{iff}\quad \exists a\in\Delta^{\mathcal{I}},\ \mathcal{I}(a:C)=\vt,\\
\mathcal{I}\models\langle\_,\_\rangle:R
 &\quad\text{iff}\quad \exists a,b\in\Delta^{\mathcal{I}},\ \mathcal{I}(\langle a,b\rangle:R)=\vt.
\end{align*}
Each assertion has its own existential scope. In particular, the two roles in $\langle\_,\_\rangle:r\sqcap s$ must hold of the same pair. Different assertions need not use distinct witnesses. Satisfaction of sets of target assertions and entailment are defined as in Definitions~\ref{def:ALC-Sat} and~\ref{def:ALC-Ent}, with these two additional kinds of assertion.
\end{definition}

An existential premise can be implemented by a fresh individual name, or a fresh pair of names, interpreted as its witnesses. Every model of the existential premises extends to such a named interpretation, and forgetting the fresh names gives the converse. This preserves consequences in the original vocabulary. The same replacement is not valid in a query: $\_:C$ may hold while a fresh name denotes an element outside $C$. Accordingly, whenever blank nodes occur on the right of an entailment below, their meaning is existential query entailment. The target consists of classical $\mathcal{ALCH}^{+}(\cap)$ together with these explicitly interpreted assertions.

\begin{definition}\label{def:tau-mu}
We define two translations $\cdot^{\tau}$ and $\cdot^{\mu}$ in tandem. Their action on atomic concepts and roles is given by the fresh copies above. Their action on arbitrary complex concepts is defined recursively:
\begin{align*}
\top^{\tau}&=\top^{\mu}=\top,\\
\bot^{\tau}&=\bot,\qquad \bot^{\mu}=\top,\\
(C\sqcap D)^{\tau}&=C^{\tau}\sqcap D^{\tau},\\
(C\sqcap D)^{\mu}&=C^{\mu}\sqcap D^{\mu},\\
(C\sqcup D)^{\tau}&=(C^{\tau}\sqcup D^{\tau})\sqcap(C^{\mu}\sqcap D^{\mu}),\\
(C\sqcup D)^{\mu}&=C^{\mu}\sqcap D^{\mu},\\
(\exists r.C)^{\tau}&=\exists r^{\tau}.C^{\tau},\\
(\exists r.C)^{\mu}&=\exists r^{\mu}.C^{\mu},\\
(\forall r.C)^{\tau}&=(\forall r^{\tau}.C^{\tau})\sqcap(\exists r^{\mu}.C^{\mu}),\\
(\forall r.C)^{\mu}&=\exists r^{\mu}.C^{\mu}.
\end{align*}
For intersection roles, $(R\sqcap S)^{\tau}=R^{\tau}\sqcap S^{\tau}$ and $(R\sqcap S)^{\mu}=R^{\mu}\sqcap S^{\mu}$. On sentences the translations are:
\begin{align*}
(a:C)^{\tau}&=a:C^{\tau}, & (a:C)^{\mu}&=a:C^{\mu},\\
(\langle a,b\rangle:R)^{\tau}&=\langle a,b\rangle:R^{\tau},\\
(\langle a,b\rangle:R)^{\mu}&=\langle a,b\rangle:R^{\mu},\\
(C\sqsubseteq D)^{\tau}&=C^{\tau}\sqsubseteq D^{\tau},\\
(C\sqsubseteq D)^{\mu}&=\_:C^{\mu}\sqcap D^{\mu},\\
(C\equiv D)^{\tau}&=C^{\tau}\equiv D^{\tau},\\
(C\equiv D)^{\mu}&=\_:C^{\mu}\sqcap D^{\mu},\\
(r\sqsubseteq s)^{\tau}&=r^{\tau}\sqsubseteq s^{\tau},\\
(r\sqsubseteq s)^{\mu}&=\langle\_,\_\rangle:r^{\mu}\sqcap s^{\mu}.
\end{align*}
Here $C,D$ range over all positive concepts, $r,s$ over atomic roles, and $R,S$ over intersection roles. In particular, the inclusion and assertion clauses apply to complex concepts, not only to atomic ones. For a set $\mathcal{S}$ of sentences, write $\mathcal{S}^{\tau}=\{\psi^{\tau}\mid\psi\in\mathcal{S}\}$ and $\mathcal{S}^{\mu}=\{\psi^{\mu}\mid\psi\in\mathcal{S}\}$.
\end{definition}

The two extra meaningfulness conditions in the truth translation are essential. Weak Kleene disjunction is indeterminate when either operand is indeterminate. A restricted universal is true only if, at the object where it is evaluated, some role--concept pair is jointly meaningful. This witness may have values $\langle\vf,\vf\rangle$; it need not be a true role edge.

\begin{definition}[Coherence axioms]\label{def:coherence}
For a finite set $\mathcal{S}$ of weak Kleene sentences, let $\mathcal{B}_{\mathcal{S}}$ contain $C^{\tau}\sqsubseteq C^{\mu}$ for every atomic concept $C$ occurring in $\mathcal{S}$, and $r^{\tau}\sqsubseteq r^{\mu}$ for every atomic role $r$ occurring in $\mathcal{S}$. For an entailment problem $\mathcal{K}=\langle\mathcal{T},\mathcal{A}\rangle$, $\varphi$, we use $\mathcal{B}_{\mathcal{T}\cup\mathcal{A}\cup\{\varphi\}}$. The query contributes coherence axioms for its symbols but is not added as a premise.
\end{definition}

All three components, the truth premises, the meaningfulness premises, and the coherence axioms, belong to the same classical theory. We next establish the model correspondence that justifies this construction.

\begin{definition}\label{def:classical-encoding}
Given a weak Kleene interpretation $\mathcal{I}$, define a classical interpretation $\mathcal{I}'$ with the same domain and individual assignments. For each atomic concept $C$, set
\begin{align*}
\mathcal{I}'(a:C^{\tau})=\vt
 &\quad\text{iff}\quad\mathcal{I}(a:C)=\vt,\\
\mathcal{I}'(a:C^{\mu})=\vt
 &\quad\text{iff}\quad\mathcal{I}(a:C)\neq\vu,
\end{align*}
and give the copied atomic roles the analogous values at each pair of domain elements. The values on all compound concepts and sentences are then determined by the classical interpretation conditions.

Conversely, given a classical interpretation $\mathcal{J}$, define a weak Kleene interpretation $\mathcal{J}'$ on the same domain and individuals. For atomic concepts put
\[
\mathcal{J}'(a:C)=\begin{cases}
\vt&\text{if }\mathcal{J}(a:C^{\tau})=\vt,\\
\vf&\text{if }\mathcal{J}(a:C^{\tau})=\vf\text{ and }\mathcal{J}(a:C^{\mu})=\vt,\\
\vu&\text{if }\mathcal{J}(a:C^{\tau})=\vf\text{ and }\mathcal{J}(a:C^{\mu})=\vf,
\end{cases}
\]
and decode atomic roles in the same way. The weak Kleene interpretation conditions determine all remaining values.
\end{definition}

Every encoding $\mathcal{I}'$ satisfies the coherence axioms. Where those axioms hold, the decoder gives $\vt$ exactly on the truth predicate and a value other than $\vu$ exactly on the meaningfulness predicate. Encoding and then decoding a weak Kleene interpretation therefore recovers its atomic assignments, domain, and individuals. Coherence is needed only for the symbols occurring in the sentences under consideration.

\begin{lemma}\label{lem:concept-correspondence}
Let $C$ be any positive concept, and suppose $\mathcal{J}$ satisfies the coherence axioms for its occurring symbols. Then, for every domain element $a$,
\begin{align*}
\mathcal{J}(a:C^{\tau})=\vt
 &\quad\text{iff}\quad\mathcal{J}'(a:C)=\vt,\\
\mathcal{J}(a:C^{\mu})=\vt
 &\quad\text{iff}\quad\mathcal{J}'(a:C)\neq\vu.
\end{align*}
The analogous equivalences hold for intersection roles. They also hold between any weak Kleene interpretation $\mathcal{I}$ and its classical encoding $\mathcal{I}'$.
\end{lemma}

\begin{proof}
We prove the two concept equivalences simultaneously by induction on the complexity of $C$.

\noindent$\bullet$ For an atomic concept the result is the decoding definition together with its coherence axiom. For $\top$ both translations are true. For $\bot$ the truth translation is false and the meaningfulness translation is true, as required by the fixed weak Kleene value $\vf$.

\noindent$\bullet$ A weak Kleene conjunction is true exactly when both operands are true, and it is meaningful exactly when both operands are meaningful. A weak Kleene disjunction is true exactly when at least one operand is true and both operands are meaningful; it too is meaningful exactly when both operands are meaningful. Substitution of the induction hypotheses gives the conjunction and disjunction clauses of Definition~\ref{def:tau-mu}.

\noindent$\bullet$ For $\exists r.C$, let
\[
X=\{\langle\mathcal{J}'(\langle a,b\rangle:r),\mathcal{J}'(b:C)\rangle\mid b\in\Delta^{\mathcal{J}}\}.
\]
By the definition of $\dot{\exists}$, its value on $X$ is $\vt$ exactly when $\langle\vt,\vt\rangle\in X$. The role correspondence and the induction hypothesis identify this condition with $\mathcal{J}(a:\exists r^{\tau}.C^{\tau})=\vt$. The value is not $\vu$ exactly when some pair in $X$ has neither coordinate equal to $\vu$. This is equivalent to $\mathcal{J}(a:\exists r^{\mu}.C^{\mu})=\vt$. A true pair is itself jointly meaningful, so the truth case is included in this characterization of meaningfulness.

\noindent$\bullet$ For $\forall r.C$, use the same set $X$. The definition of $\dot{\forall}$ requires both the absence of $\langle\vt,\vf\rangle$ and $\langle\vt,\vu\rangle$ and the existence of a jointly meaningful pair. Absence of those two pairs says that every $b$ with $\mathcal{J}'(\langle a,b\rangle:r)=\vt$ has $\mathcal{J}'(b:C)=\vt$. It therefore translates to $a:\forall r^{\tau}.C^{\tau}$. The witness condition translates to $a:\exists r^{\mu}.C^{\mu}$. Their conjunction gives the truth translation. Finally, $\dot{\forall}(X)\neq\vu$ exactly when that witness condition holds, giving the meaningfulness translation.

For intersection roles, induction on the role expression uses the conjunction truth table in exactly the same way as the concept-conjunction case. The encoding direction follows by applying the proved equivalences to $\mathcal{I}'$ and observing that its decoding recovers $\mathcal{I}$.
\end{proof}

\begin{lemma}\label{lem:sentence-correspondence}
Suppose $\mathcal{J}$ satisfies the coherence axioms for a weak Kleene sentence $\psi$. Then
\[
\mathcal{J}'(\psi)=\vt\quad\text{iff}\quad
\mathcal{J}\models\psi^{\tau}\text{ and }\mathcal{J}\models\psi^{\mu}.
\]
The same equivalence holds for $\mathcal{I}$ and its classical encoding $\mathcal{I}'$.
\end{lemma}

\begin{proof}
For concept and role assertions, Lemma~\ref{lem:concept-correspondence} applies directly: truth implies meaningfulness, so requiring both translations is equivalent to requiring truth.

For $C\sqsubseteq D$, the set of input pairs to $\dot{\forall}$ is
\[
\{\langle\mathcal{J}'(a:C),\mathcal{J}'(a:D)\rangle\mid a\in\Delta^{\mathcal{J}}\}.
\]
Its universal value is $\vt$ exactly when every true antecedent has a true consequent and some object gives a jointly meaningful pair. Lemma~\ref{lem:concept-correspondence} translates these conditions to $C^{\tau}\sqsubseteq D^{\tau}$ and $\_:C^{\mu}\sqcap D^{\mu}$, respectively.

For $C\equiv D$, weak Kleene conjunction is true exactly when both directed inclusions are true. Their truth conditions combine into $C^{\tau}\equiv D^{\tau}$. Their witness conditions coincide, because an object satisfies $C^{\mu}\sqcap D^{\mu}$ exactly when it satisfies $D^{\mu}\sqcap C^{\mu}$.

For $r\sqsubseteq s$, the input pairs are indexed by $\Delta^{\mathcal{J}}\times\Delta^{\mathcal{J}}$. The same reasoning gives $r^{\tau}\sqsubseteq s^{\tau}$ and the existence of a pair satisfying both $r^{\mu}$ and $s^{\mu}$, which is precisely $\langle\_,\_\rangle:r^{\mu}\sqcap s^{\mu}$. The encoding assertion follows from Lemma~\ref{lem:concept-correspondence} in the same cases.
\end{proof}

\begin{lemma}\label{lem:alh-wk-to-alh}
Let $\mathcal{K}=\langle\mathcal{T},\mathcal{A}\rangle$ be a weak Kleene knowledge base and $\varphi$ a weak Kleene sentence. Form the single classical knowledge base
\[
\left\langle
\mathcal{T}^{\tau}\cup\mathcal{B}_{\mathcal{T}\cup\mathcal{A}\cup\{\varphi\}},\;
\mathcal{A}^{\tau}\cup\mathcal{T}^{\mu}\cup\mathcal{A}^{\mu}
\right\rangle.
\]
Then $\mathcal{K}\models_{\mathcal{ALCH}^{+}(\cap)_{\mathsf{wK}}}\varphi$ iff this classical knowledge base entails both $\varphi^{\tau}$ and $\varphi^{\mu}$, with blank-node conclusions interpreted as in Definition~\ref{def:blank-nodes}.
\end{lemma}

\begin{proof}
\noindent$\bullet$ For left-to-right, let $\mathcal{J}$ be any model of the displayed classical knowledge base. Its coherence axioms cover all the premises and the query. By Lemma~\ref{lem:sentence-correspondence}, its decoding $\mathcal{J}'$ satisfies every sentence in $\mathcal{T}\cup\mathcal{A}$. The assumed weak Kleene entailment gives $\mathcal{J}'(\varphi)=\vt$. Applying that lemma to the query yields both $\mathcal{J}\models\varphi^{\tau}$ and $\mathcal{J}\models\varphi^{\mu}$.

\noindent$\bullet$ For right-to-left, let $\mathcal{I}$ be any weak Kleene model of $\mathcal{K}$. Its encoding $\mathcal{I}'$ satisfies all the coherence axioms. Lemma~\ref{lem:sentence-correspondence} shows that it satisfies both translations of every premise, so it is a model of the displayed classical knowledge base. The two assumed target entailments give $\mathcal{I}'\models\varphi^{\tau}$ and $\mathcal{I}'\models\varphi^{\mu}$, whence $\mathcal{I}(\varphi)=\vt$.

No satisfiability assumption is used. For example, if $a:\bot$ is a premise, its truth translation is also $a:\bot$, so both source and target have no model. Their entailments are therefore vacuous in the same way. This would not be ensured by checking the two translations against separate knowledge bases.
\end{proof}

\subsection{Embedding $\mathcal{ALCH}_{\mathsf{AC}}$ in $\mathcal{ALCH}^+(\cap)_{\mathsf{wK}}$}

Our interest in $\mathcal{ALCH}^{+}(\cap)_{\mathsf{wK}}$ is functional: its translation provides the intermediate step for $\mathcal{ALCH}_{\mathsf{AC}}$. Following the bilateral analysis in \cite{DBLP:journals/logcom/Ferguson16}, we use $\cdot^{\pi}$ and $\cdot^{\nu}$ to record the verification and falsification coordinates of concepts. Both are positive weak Kleene concepts. The sentence translation needed for entailment is $\cdot^{\pi}$, because satisfaction in $\mathsf{AC}$ requires its first coordinate to be $\vt$.

\begin{definition}\label{def:pi-nu}
For each atomic concept $C$ and role $r$, introduce distinct weak Kleene copies $C^{\pi},C^{\nu}$ and $r^{\pi},r^{\nu}$. Let $C_{\top}^{\nu}$ be a further fresh concept name, and use $C_{\bot}^{\pi}$ as a second notation for that same name. This identification preserves the common value $v$ in the two constant clauses of Definition~\ref{def:Bilateral-Interp}. For arbitrary concepts, define recursively:
\begin{align*}
\top^{\pi}&=\top, &\top^{\nu}&=C_{\top}^{\nu},\\
\bot^{\pi}&=C_{\bot}^{\pi}, &\bot^{\nu}&=\top,\\
(\neg C)^{\pi}&=C^{\nu}, &(\neg C)^{\nu}&=C^{\pi},\\
(C\sqcap D)^{\pi}&=C^{\pi}\sqcap D^{\pi},
 &(C\sqcap D)^{\nu}&=C^{\nu}\sqcup D^{\nu},\\
(C\sqcup D)^{\pi}&=C^{\pi}\sqcup D^{\pi},
 &(C\sqcup D)^{\nu}&=C^{\nu}\sqcap D^{\nu},\\
(\exists r.C)^{\pi}&=\exists r^{\pi}.C^{\pi},
 &(\exists r.C)^{\nu}&=\forall r^{\pi}.C^{\nu},\\
(\forall r.C)^{\pi}&=\forall r^{\pi}.C^{\pi},
 &(\forall r.C)^{\nu}&=\exists r^{\pi}.C^{\nu}.
\end{align*}
On assertions, put
\begin{align*}
(a:C)^{\pi}&=a:C^{\pi}, &(a:C)^{\nu}&=a:C^{\nu},\\
(\langle a,b\rangle:r)^{\pi}&=\langle a,b\rangle:r^{\pi},\\
(\langle a,b\rangle:r)^{\nu}&=\langle a,b\rangle:r^{\nu}.
\end{align*}
Extend $\cdot^{\pi}$ to the remaining sentences by
\begin{align*}
(C\sqsubseteq D)^{\pi}&=C^{\pi}\sqsubseteq D^{\pi},\\
(C\equiv D)^{\pi}&=C^{\pi}\equiv D^{\pi},\\
(r\sqsubseteq s)^{\pi}&=r^{\pi}\sqsubseteq s^{\pi}.
\end{align*}
No negative translation of inclusions is needed for the entailment theorem. The translation $\cdot^{\nu}$ on concepts already supplies the negative coordinates required by arbitrary nested concept negations.
\end{definition}

We now define the interpretations associated with these translations. Write $\pi_0$ and $\pi_1$ for the first and second projections on $\mathcal{V}_3^2$, respectively.

\begin{definition}\label{def:coordinate-interpretations}
For an $\mathsf{AC}$ interpretation $\mathcal{I}$, define a weak Kleene interpretation $\mathcal{I}^{\circ}$ on the copied signature, preserving its domain and individual assignments. For atomic concepts and roles put
\begin{align*}
\mathcal{I}^{\circ}(a:C^{\pi})&=\pi_0(\mathcal{I}(a:C)),\\
\mathcal{I}^{\circ}(a:C^{\nu})&=\pi_1(\mathcal{I}(a:C)),\\
\mathcal{I}^{\circ}(\langle a,b\rangle:r^{\pi})&=\pi_0(\mathcal{I}(\langle a,b\rangle:r)),\\
\mathcal{I}^{\circ}(\langle a,b\rangle:r^{\nu})&=\pi_1(\mathcal{I}(\langle a,b\rangle:r)).
\end{align*}
The additional constant name, written in its two notations, is assigned
\begin{align*}
\mathcal{I}^{\circ}(a:C_{\top}^{\nu})&=\pi_1(\mathcal{I}(a:\top)),\\
\mathcal{I}^{\circ}(a:C_{\bot}^{\pi})&=\pi_0(\mathcal{I}(a:\bot)).
\end{align*}
Conversely, for a weak Kleene interpretation $\mathcal{I}$ on the copied signature, define $\mathcal{I}^{\bullet}$ on the original signature by
\begin{align*}
\mathcal{I}^{\bullet}(a:C)
 &=\langle\mathcal{I}(a:C^{\pi}),\mathcal{I}(a:C^{\nu})\rangle,\\
\mathcal{I}^{\bullet}(\langle a,b\rangle:r)
 &=\langle\mathcal{I}(\langle a,b\rangle:r^{\pi}),\mathcal{I}(\langle a,b\rangle:r^{\nu})\rangle
\end{align*}
for atomic concepts and roles, and by
\begin{align*}
\mathcal{I}^{\bullet}(a:\top)&=\langle\vt,\mathcal{I}(a:C_{\top}^{\nu})\rangle,\\
\mathcal{I}^{\bullet}(a:\bot)&=\langle\mathcal{I}(a:C_{\bot}^{\pi}),\vt\rangle
\end{align*}
for the constants. Preserve the domain and individual assignments, and extend by the $\mathsf{AC}$ interpretation conditions.
\end{definition}

\begin{lemma}\label{lem:ALCH-AC-interpretations}
Let $\mathcal{I}$ be an $\mathsf{AC}$ interpretation. For every concept $C$ and domain element $a$,
\begin{align*}
\mathcal{I}^{\circ}(a:C^{\pi})&=\pi_0(\mathcal{I}(a:C)),\\
\mathcal{I}^{\circ}(a:C^{\nu})&=\pi_1(\mathcal{I}(a:C)).
\end{align*}
For every $\mathcal{ALCH}_{\mathsf{AC}}$ sentence $\varphi$,
\[
\mathcal{I}^{\circ}(\varphi^{\pi})=\pi_0(\mathcal{I}(\varphi)).
\]
\end{lemma}

\begin{proof}
We prove the concept assertions together by induction on $C$. Atomic concepts follow from Definition~\ref{def:coordinate-interpretations}. The constant clauses hold because the single name denoted by $C_{\top}^{\nu}$ and $C_{\bot}^{\pi}$ records their common unspecified value $v$. The two assignments to this name in Definition~\ref{def:coordinate-interpretations} agree by the original constant clause. Negation exchanges the coordinates, so the two induction hypotheses exchange their roles. Conjunction uses $\dot{\wedge}$ on the first coordinates and $\dot{\vee}$ on the second; disjunction uses the opposite pair. These are precisely the corresponding clauses of Definition~\ref{def:pi-nu}.

For an existential restriction, let
\[
X=\{\langle\mathcal{I}(\langle a,b\rangle:r),\mathcal{I}(b:C)\rangle\mid b\in\Delta^{\mathcal{I}}\}.
\]
By Definition~\ref{def:AC-quantifiers}, its two coordinates are $\dot{\exists}(\pi^{0,0}[X])$ and $\dot{\forall}(\pi^{0,1}[X])$. The definition of $\mathcal{I}^{\circ}$ and the induction hypotheses identify these input sets with
\begin{align*}
&\{\langle\mathcal{I}^{\circ}(\langle a,b\rangle:r^{\pi}),\mathcal{I}^{\circ}(b:C^{\pi})\rangle
       \mid b\in\Delta^{\mathcal{I}}\},\\
&\{\langle\mathcal{I}^{\circ}(\langle a,b\rangle:r^{\pi}),\mathcal{I}^{\circ}(b:C^{\nu})\rangle
       \mid b\in\Delta^{\mathcal{I}}\},
\end{align*}
respectively. Hence the two values are the weak Kleene values of $a:\exists r^{\pi}.C^{\pi}$ and $a:\forall r^{\pi}.C^{\nu}$. For a universal restriction, Definition~\ref{def:AC-quantifiers} exchanges $\dot{\exists}$ and $\dot{\forall}$, giving the universal translation in exactly the same manner. This proves the concept part for all nesting depths.

For sentences, concept assertions follow from the proved first-coordinate equality and role assertions from the copied role assignment. The first coordinate of $C\sqsubseteq D$ is $\dot{\forall}$ applied to the pairs of first coordinates of $C$ and $D$; the concept equalities therefore identify it with $\mathcal{I}^{\circ}(C^{\pi}\sqsubseteq D^{\pi})$. Equivalence uses $\dot{\wedge}$ on the first coordinates of its two directed inclusions. Role inclusion has the same argument as concept inclusion, with pairs indexed by the product domain. This proves the sentence equality.
\end{proof}

\begin{corollary}\label{cor:AC-to-wK}
For every $\mathsf{AC}$ interpretation $\mathcal{I}$ and source sentence $\varphi$,
\[
\mathcal{I}\models_{\mathcal{ALCH}_{\mathsf{AC}}}\varphi
\quad\text{iff}\quad
\mathcal{I}^{\circ}\models_{\mathcal{ALCH}^{+}(\cap)_{\mathsf{wK}}}\varphi^{\pi}.
\]
\end{corollary}
\begin{proof}
Both sides require the corresponding value in Lemma~\ref{lem:ALCH-AC-interpretations} to be $\vt$.
\end{proof}

\begin{corollary}\label{cor:wk-to-AC}
For every weak Kleene interpretation $\mathcal{I}$ on the copied signature and every source $\mathcal{ALCH}_{\mathsf{AC}}$ sentence $\varphi$,
\[
\mathcal{I}^{\bullet}\models_{\mathcal{ALCH}_{\mathsf{AC}}}\varphi
\quad\text{iff}\quad
\mathcal{I}\models_{\mathcal{ALCH}^{+}(\cap)_{\mathsf{wK}}}\varphi^{\pi}.
\]
\end{corollary}
\begin{proof}
The interpretation $(\mathcal{I}^{\bullet})^{\circ}$ recovers $\mathcal{I}$ on every copied atomic concept and role, on the additional constant name, and on the domain and individuals. It therefore agrees with $\mathcal{I}$ on all intermediate sentences. Apply Corollary~\ref{cor:AC-to-wK} to $\mathcal{I}^{\bullet}$.
\end{proof}

\begin{lemma}\label{lem:alch-ac-to-alh-wk}
The translation $\cdot^{\pi}$ faithfully translates $\mathcal{ALCH}_{\mathsf{AC}}$ entailment into $\mathcal{ALCH}^{+}(\cap)_{\mathsf{wK}}$ entailment:
\begin{align*}
\langle\mathcal{T},\mathcal{A}\rangle
 &\models_{\mathcal{ALCH}_{\mathsf{AC}}}\varphi
\quad\text{iff}\\
\langle\mathcal{T}^{\pi},\mathcal{A}^{\pi}\rangle
 &\models_{\mathcal{ALCH}^{+}(\cap)_{\mathsf{wK}}}\varphi^{\pi}.
\end{align*}
\end{lemma}
\begin{proof}
Suppose the source entailment holds, and let $\mathcal{I}$ be a weak Kleene model of the translated premises. Corollary~\ref{cor:wk-to-AC} makes $\mathcal{I}^{\bullet}$ an $\mathsf{AC}$ model of the original premises. It satisfies $\varphi$ by the assumption, so the same corollary gives $\mathcal{I}\models\varphi^{\pi}$. Conversely, any $\mathsf{AC}$ model of the original premises projects by Corollary~\ref{cor:AC-to-wK} to a weak Kleene model of the translated premises. Apply the intermediate entailment and that corollary again to obtain satisfaction of the original query.
\end{proof}

We can now state the main theorem. Using the composition notation, $\varphi^{\tau\circ\pi}$ means $(\varphi^{\pi})^{\tau}$, and similarly for $\mu\circ\pi$.

\begin{theorem}\label{thm:faithful-ac}
Let $\mathcal{K}=\langle\mathcal{T},\mathcal{A}\rangle$ be an $\mathcal{ALCH}_{\mathsf{AC}}$ knowledge base and $\varphi$ a source sentence. Add the coherence axioms for every symbol occurring in $\mathcal{T}^{\pi}\cup\mathcal{A}^{\pi}\cup\{\varphi^{\pi}\}$ and form the single classical knowledge base with TBox
\[
\mathcal{T}^{\tau\circ\pi}\cup
\mathcal{B}_{\mathcal{T}^{\pi}\cup\mathcal{A}^{\pi}\cup\{\varphi^{\pi}\}}
\]
and ABox
\[
\mathcal{A}^{\tau\circ\pi}\cup
\mathcal{T}^{\mu\circ\pi}\cup\mathcal{A}^{\mu\circ\pi}.
\]
Then $\mathcal{K}\models_{\mathcal{ALCH}_{\mathsf{AC}}}\varphi$ iff this classical knowledge base entails both $\varphi^{\tau\circ\pi}$ and $\varphi^{\mu\circ\pi}$. The target uses the existential interpretation of blank-node assertions from Definition~\ref{def:blank-nodes}.
\end{theorem}

\begin{proof}
By Lemma~\ref{lem:alch-ac-to-alh-wk}, source entailment is equivalent to the weak Kleene entailment of $\varphi^{\pi}$ from $\langle\mathcal{T}^{\pi},\mathcal{A}^{\pi}\rangle$. Apply Lemma~\ref{lem:alh-wk-to-alh} to that entailment. Its truth and meaningfulness premises, together with its coherence axioms, are exactly the TBox and ABox displayed above. Its two conclusions are exactly $\varphi^{\tau\circ\pi}$ and $\varphi^{\mu\circ\pi}$.

Equivalently, the proof gives countermodel conversions in both directions. An $\mathsf{AC}$ countermodel projects to a weak Kleene countermodel and then encodes as a classical model of all the displayed premises in which at least one translated conclusion fails. Conversely, a classical model of those premises failing a translated conclusion decodes to a weak Kleene countermodel, which lifts to an $\mathsf{AC}$ countermodel. Domains and individual assignments are preserved; neither finiteness of the domain nor satisfiability of the knowledge base is assumed.
\end{proof}

The theorem handles arbitrary complex concepts in assertions and inclusions, including nested negations and role restrictions. It establishes a reduction to classical reasoning with explicitly existential queries. When an existential assertion is a premise, fresh named witnesses give the ordinary ABox presentation described above. When it is a query, its existential meaning must be retained. In particular, the theorem does not identify this task with ordinary named-assertion entailment alone.

The additional restriction to $\mathcal{ALCH}_{\mathsf{S}^*}$ requires a separate correspondence between its admissible interpretations and additional classical axioms. The theorem proved here concerns $\mathcal{ALCH}_{\mathsf{AC}}$; no entailment theorem for the five-valued restriction is used in its proof.

\subsection{Computational Complexity}

The translations are defined by structural recursion, and an entailment problem with a finite TBox and ABox produces finitely many translated sentences and coherence axioms. This remains true when the ambient signature is infinite, since coherence axioms are required only for symbols in the premises and the query.

Finiteness does not by itself establish a linear bound for the expanded syntax trees. For example, the truth translation of $\forall r.C$ contains both $C^{\tau}$ and $C^{\mu}$. Representing repeated subexpressions by shared definitions may avoid duplication, but a size bound and a proof that those definitions preserve the relevant models would be needed for such an implementation. Moreover, the target includes existential queries, so complexity results for ordinary axiom entailment cannot be applied without a further argument. We therefore do not infer tight \textsf{PSPACE} or \textsf{EXPTIME} bounds from the faithful translation. Matching lower bounds would also require a reduction in the opposite direction or an independently established hard fragment.

\section{Discussion}

One can immediately note relationships with the work of Maier \emph{et al.} in \cite{DBLP:journals/semweb/MaierMH13} on non-classical versions of OWL in the form of the four-valued description logic $\mathcal{SROIQ}_{4}$. The expressivity in the description logics in the present paper has increased that of \cite{DBLP:conf/dlog/Ferguson21}, and porting the present intuitions to such a language is an attractive goal. Considering such a transition of the work to a more expressive paradigm like $\mathcal{SROIQ}$ prompts us to acknowledge some extra-logical design decisions that go into the definition of a non-classical description logic.

Note that in \cite{DBLP:journals/semweb/MaierMH13}, as in the present work, interpreting the falsification of an assertion of a role restriction $a:\forall r.C$ is essentially identified with the verification of the role restriction $a:\exists r.\neg C$. The details of \cite{DBLP:journals/semweb/MaierMH13} make explicit what we have assumed implicit in the present work, namely, that the interpretation essentially rests on a more primitive interpretation of the verification and falsification of \emph{the conditional}. Both the present and related work (not only \cite{DBLP:journals/semweb/MaierMH13} but also that of Norihiro Kamide in \cite{DBLP:conf/kes/Kamide11}) implicitly follow the lead of David Nelson's analysis of \emph{constructible falsity} in \emph{e.g.} \cite{DBLP:journals/jsyml/AlmukdadN84} through an identification of a negated conditional $\neg(\varphi\rightarrow\psi)$ with the conjunction of the antecedent $\varphi$ and the negation of the consequent $\neg\psi$. This observation presses two related degrees of freedom that are independent of the mere choice of truth values, \emph{i.e.}, the matter of the semantics of the conditional and those of the negated conditional.

In the first case, the related context of the four-valued $\mathcal{ALC}4$ introduced by Ma \emph{et al.} in \cite{DBLP:conf/esws/MaHL07} explicitly acknowledges several options for the interpretation of the conditional to which role restrictions implicitly appeal. Such a context distinguishes between the \emph{material} implication and the \emph{intensional} conditional, each of which will induce a different semantics for role restrictions. It is \emph{e.g.} a different matter that $r\sqsubseteq s$ should hold in virtue of the fact that all pairs locally known to satisfy $r$ also satisfy $s$ than that $r\sqsubseteq s$ should encode a deeper relationship between the roles $r$ and $s$ beyond a contingent fact about known data.

The present bilateral context also underscores the nontriviality of assessing the negation of a conditional, an account which must be solved for negated role restrictions. As Heinrich Wansing stresses in \cite{DBLP:conf/aiml/Wansing10}, the Nelsonian account is not inevitable and there are in fact several \emph{a priori} natural options. Wansing, for example, prefers an account in which $\neg(\varphi\rightarrow\psi)$ is identified with $\varphi\rightarrow\neg\psi$. Just as choice of conditional is an independent matter from the choice of truth values, so, too, is choice of negated conditional. 

Our bilateral context stresses that both dimensions be taken seriously. As expressivity of the language grows when providing versions of OWL in line with the present approach, such design decisions become applicable to further cases. 


This observation has consequences also for the role that the present work can play in enhancing the significant increase in accuracy won by adopting bilateral reasoning by Allen \emph{et al.} in \cite{DBLP:conf/nesy/AllenCFIG25,allenetal2026}. These papers show improved accuracy in distilling semantic information from a large language model but stop short at the atomic level. Extending such an analysis of factuality from the atomic level to the full language of $\mathcal{ALC}$ or $\mathcal{SROIQ}$ is thus important for two related reasons. As a natural extension to the scope of \cite{DBLP:conf/nesy/AllenCFIG25,allenetal2026}, empirically testing whether bilateral interpretations of \emph{e.g.} role restrictions should improve over classical description logic methods is worth pursuing. But from the other direction, such experiments should guide the design of description logics in the vicinity of $\mathcal{ALC}_{AC}$ as we confirm which design decisions concerning verification and falsification of conditionals are best suited to particular use-cases.

\bibliography{refs-revised}

\end{document}
